% Part: conditional-logics
% Chapter: minimal-change-semantics
% Section: true-false

\documentclass[../../../include/open-logic-section]{subfiles}

\begin{document}

\olfileid{con}{min}{tf}

\olsection{ప్రతివాస్తవాల సత్యం, అసత్యం}

\olref{fig:true}లో చూపినట్లుగా,
అత్యంత సమీప $!A$-లోకాలన్నీ
$!B$-లోకాలైతే ప్రతివాస్తవం
$!A \cif !B$ (శూన్యసత్యం
కాని విధంగా) సత్యం.
\begin{figure}
\begin{center}
\begin{tikzpicture}[scale=.7]\small
  \spheresystem{5}
  \draw (0,0) node {$w$};
  \propositionintersect{0}{4}{40}{3.5}
  {\tikzset{proposition/.append={smooth,tension=1.4}}
  \proposition[shift={(1.6,-1)}]{90}{1}{30}{4}}
  \path (1.5,3) node[below] {$\formula{B}$};
  \path (3,0) node {$\formula{A}$};
\end{tikzpicture}
\caption{శూన్యసత్యం కాని సత్య ప్రతివాస్తవం}
\ollabel{fig:true}
\end{center}
\end{figure}
~$w$ చుట్టూ ఉన్న గోళ
వ్యవస్థలో అసలు ఏ $!A$-లోకాన్నీ
కలిగిన గోళం లేకపోయినా, ఆ
ప్రతివాస్తవం~$w$ వద్ద సత్యమే.
ఈ సందర్భంలో అది \emph{శూన్యసత్యం}
(\olref{fig:vacuous} చూడండి).
\begin{figure}
\begin{center}
\begin{tikzpicture}[scale=.7]\small
  \spheresystem{5}
  \draw (0,0) node {$w$};
  \proposition{0}{6.5}{40}{3.5}
  {\tikzset{proposition/.append={smooth,tension=1.4}}
  \proposition[shift={(1.2,-1)}]{90}{1}{30}{4}}
  \path (1.5,3) node[below] {$\formula{B}$};
  \spherepos{0}{7}{node {$\formula{A}$}}
\end{tikzpicture}
\caption{శూన్యసత్య ప్రతివాస్తవం}
\ollabel{fig:vacuous}
\end{center}
\end{figure}

ప్రతివాస్తవం రెండు విధాలుగా
అసత్యమవుతుంది. ఒక విధం:
అత్యంత సమీప $!A$-లోకాలన్నీ
$!B$-లోకాలు కావు, కానీ
కొన్ని మాత్రం అటువంటివే.
ఈ సందర్భంలో $!A \cif
\lnot !B$ కూడా అసత్యమే
(\olref{fig:false} చూడండి).
\begin{figure}
\begin{center}
\begin{tikzpicture}[scale=.7]\small
  \spheresystem{5}
  \draw (0,0) node {$w$};
  \propositionintersect{0}{4}{40}{3.5}
  {\tikzset{proposition/.append={smooth,tension=1.4}}
  \proposition[shift={(1.6,0)}]{90}{1}{30}{3}}
  \path (1.5,3) node[below] {$\formula{B}$};
  \path (3,0) node {$\formula{A}$};
\end{tikzpicture}
\caption{అసత్య ప్రతివాస్తవం; వ్యతిరేకమూ అసత్యం}
\ollabel{fig:false}
\end{center}
\end{figure}
అత్యంత సమీప $!A$-లోకాలతో
$!B$-లోకాలకు ఏ అతివ్యాప్తీ
లేకపోతే, $!A \cif !B$
అసత్యం. కానీ ఈ సందర్భంలో
అత్యంత సమీప $!A$-లోకాలన్నీ
$\lnot !B$-లోకాలు; అందువల్ల
$!A \cif \lnot !B$ సత్యం
(\olref{fig:false-opposite} చూడండి).
\begin{figure}
\begin{center}
\begin{tikzpicture}[scale=.7]\small
  \spheresystem{5}
  \draw (0,0) node {$w$};
  \propositionintersect{0}{4}{40}{3.5}
  {\tikzset{proposition/.append={smooth,tension=1.4}}
  \proposition[shift={(-1,0)}]{90}{1}{30}{3}}
  \path (-1,3) node[below] {$\formula{B}$};
  \path (3,0) node {$\formula{A}$};
\end{tikzpicture}
\caption{అసత్య ప్రతివాస్తవం; వ్యతిరేకం సత్యం}
\ollabel{fig:false-opposite}
\end{center}
\end{figure}

కఠిన సోపాధికానికి భిన్నంగా,
ప్రతివాస్తవాలు స్థితిని బట్టి
సత్యమో అసత్యమో కావచ్చు.
\olref{fig:contingent}లోని
గోళ నమూనాను పరిశీలించండి.
~$u$కు అత్యంత సమీప
$!A$-లోకాలన్నీ $!B$-లోకాలు;
కాబట్టి $\mSat{M}{!A \cif
  !B}[u]$. కానీ~$v$కు
అత్యంత సమీపంగా ఉన్న కొన్ని
$!A$-లోకాలు $!B$-లోకాలు
కావు; కాబట్టి
$\mSat/{M}{!A \cif !B}[v]$.
\begin{figure}
\begin{center}
\begin{tikzpicture}[scale=.6]\small
  \spheresystem{7}
  \spheresystem[shift={(0:3.1)},dashed]{4}
  \propositionintersect{45}{4}{40}{4.5}
  \begin{scope}
    \clip \propositionplot{45}{4}{40}{4.5} ;
    \spherelayer[shift={(0:3.1)}]{3}
  \end{scope}
  \proposition[shift={(1.4,-.5)}]{90}{1}{35}{5}
  \path (0,0) node {$u$};
  \path (0:3.1) node {$v$};
  \spherepos{35}{9}{node {$\formula{A}$}}
  \spherepos{80}{9}{node {$\formula{B}$}}
\end{tikzpicture}
\end{center}
\caption{స్థితిని బట్టి సత్యం మారే ప్రతివాస్తవం}
\ollabel{fig:contingent}
\end{figure}
\end{document}
